\documentclass[12pt]{article}

\usepackage{sbc-template}
\usepackage{graphicx,url}
\usepackage[utf8]{inputenc}
\usepackage[brazil]{babel}
\usepackage{listings}
\usepackage{xcolor}

\lstset{
  basicstyle=\ttfamily\footnotesize,
  breaklines=true,
  frame=single,
  backgroundcolor=\color{gray!10},
  keywordstyle=\color{blue},
  commentstyle=\color{green!60!black},
  stringstyle=\color{red!70!black},
}

\sloppy

\title{AVSYS: Um Sistema Distribuído de Reserva de Passagens Aéreas\\
  com Prevenção de Double-Booking e Comunicação Assíncrona}

\author{Felipe Farias\inst{1}, Breno Moura\inst{1},
  Gabriel Aragão\inst{1}, Nicollas Matsuo\inst{1}}

\address{Curso de Ciência da Computação\\
  [Nome da Universidade] -- [Cidade] -- [Estado] -- Brasil
  \email{\{felipe.farias, breno.moura, gabriel.aragao, nicollas.matsuo\}@[universidade].br}
}

\begin{document}

\maketitle

\begin{abstract}
This paper presents AVSYS, a distributed airline reservation system designed to handle
high-concurrency scenarios while preventing double-booking anomalies. The system adopts
a microservices architecture composed of three independent services — flight catalog,
reservation management, and payment processing — orchestrated through an NGINX API
Gateway. Consistency under concurrent access is achieved through a two-layer defense:
a distributed mutex implemented with Redis atomic \texttt{SET NX EX} commands, and
optimistic row versioning in SQLite. Asynchronous inter-service communication is handled
via RabbitMQ topic exchanges with durable queues and a Dead Letter Exchange for fault
tolerance. All services are implemented with the Bun runtime and the ElysiaJS framework,
containerized with Docker Compose. Experimental validation confirms zero double-bookings
under simulated concurrent load, with an average reservation pipeline latency of
approximately 120 ms.
\end{abstract}

\begin{resumo}
Este artigo apresenta o AVSYS, um sistema distribuído de reserva de passagens aéreas
projetado para cenários de alta concorrência com prevenção de double-booking. A
arquitetura adota o padrão de microsserviços, composta por três serviços independentes
— catálogo de voos, gerenciamento de reservas e processamento de pagamentos —
orquestrados por um API Gateway NGINX. A consistência sob acesso concorrente é garantida
por uma defesa em duas camadas: um mutex distribuído com Redis (\texttt{SET NX EX}
atômico) e versionamento otimista de linhas no SQLite. A comunicação assíncrona
entre serviços é realizada via RabbitMQ com filas duráveis e Dead Letter Exchange
para tolerância a falhas. Todos os serviços são implementados com o runtime Bun e o
framework ElysiaJS, conteinerizados com Docker Compose. A validação experimental
confirma zero double-bookings sob carga concorrente simulada, com latência média do
pipeline de reserva de aproximadamente 120 ms.
\end{resumo}


% =============================================================================
\section{Introdução}
% =============================================================================

A comercialização de passagens aéreas é um dos domínios mais exigentes para sistemas
de informação: dezenas de milhares de usuários podem tentar reservar o mesmo conjunto
limitado de assentos em intervalos de segundos, especialmente em lançamentos de
promoções ou períodos de alta demanda. Nesse contexto, garantir que dois usuários não
obtenham reservas válidas para o mesmo assento — problema conhecido como
\textit{double-booking} — é um requisito de corretude crítico, não uma otimização
desejável \cite{gray92}.

Sistemas monolíticos tradicionais contornam este problema com transações de banco de
dados e bloqueios pessimistas. Contudo, a adoção crescente de arquiteturas de
microsserviços impõe novos desafios: serviços independentes possuem bancos de dados
próprios, não há transação distribuída global, e múltiplas réplicas do mesmo serviço
podem processar requisições simultaneamente sem coordenação implícita \cite{newman2015}.
O teorema CAP \cite{brewer2000} formaliza a tensão entre consistência, disponibilidade
e tolerância a particionamento — em um sistema distribuído, os três não podem ser
garantidos simultaneamente.

Este trabalho apresenta o \textbf{AVSYS} (\textit{Airline Reservation System}),
implementado como projeto final da disciplina de Sistemas Distribuídos e Programação
Paralela. Os objetivos específicos são:

\begin{enumerate}
  \item Projetar uma arquitetura de microsserviços com fronteiras de dados claras e
        comunicação assíncrona entre serviços;
  \item Implementar um mecanismo robusto de prevenção de double-booking através de
        lock distribuído Redis combinado com versionamento otimista SQLite;
  \item Demonstrar o uso de mensageria assíncrona (RabbitMQ) para desacoplar o
        processamento de pagamentos do fluxo de reserva;
  \item Containerizar o sistema completo com Docker Compose, incluindo health checks
        e redes isoladas;
  \item Desenvolver um painel de visualização em tempo real das interações entre
        serviços.
\end{enumerate}

O restante deste artigo está organizado da seguinte forma: a Seção~2 descreve a
metodologia e arquitetura do sistema; a Seção~3 apresenta os resultados e discussão;
a Seção~4 conclui com considerações finais e trabalhos futuros.


% =============================================================================
\section{Metodologia}
% =============================================================================

\subsection{Arquitetura do Sistema}

O AVSYS foi projetado segundo o padrão arquitetural de microsserviços \cite{newman2015},
no qual uma aplicação é decomposta em serviços pequenos, independentes e implantáveis
de forma autônoma. Cada microsserviço encapsula um subdomínio de negócio bem definido,
mantém seu próprio estado persistente e se comunica com os demais por meio de interfaces
padronizadas — chamadas HTTP REST síncronas para operações onde o cliente aguarda
resposta, e mensagens assíncronas via fila para efeitos colaterais entre serviços.

Essa abordagem confere ao sistema três propriedades fundamentais: (i)~\textbf{escalabilidade
horizontal independente} — serviços com alta demanda podem ter múltiplas réplicas sem
afetar os demais; (ii)~\textbf{isolamento de falhas} — a indisponibilidade de um
serviço não propaga falhas em cascata \cite{kleppmann2017}; e
(iii)~\textbf{autonomia de dados} — cada serviço é proprietário exclusivo de seu banco
de dados, eliminando a necessidade de transações distribuídas entre serviços
\cite{richardson2018}.

A borda do sistema é composta por um API Gateway (NGINX) que recebe todas as
requisições externas, roteia-as para o microsserviço correspondente com base no prefixo
de caminho da URL, e aplica políticas de \textit{rate limiting}. A comunicação
assíncrona entre serviços ocorre por meio de um barramento de eventos baseado em
RabbitMQ. A Figura~\ref{fig:arch} ilustra a topologia geral do sistema.

\begin{figure}[ht]
\centering
\begin{lstlisting}[language={}]
             [ CLIENTE ]
                  |  HTTP :80
             [ NGINX ]  -- rate limiting, least_conn
           /     |      \
  /flights/  /reservations/  /payments/
       |           |              |
 flight-catalog  reservation   payment
   :3001          service       service
  SQLite(WAL)     :3002         :3003
                 SQLite(WAL)  SQLite(WAL)
       |           |    |         |
       |       [Redis]  |     [RabbitMQ]
       |       locks/   |    airline.events
       +-------cache----+         |
                                  +-- payment.requested -->  payment-service
                                  <-- payment.succeeded/failed --
\end{lstlisting}
\caption{Topologia da arquitetura AVSYS: três microsserviços orquestrados por NGINX,
  com Redis para locks/cache e RabbitMQ para comunicação assíncrona.}
\label{fig:arch}
\end{figure}

\subsection{Tecnologias Utilizadas}

\subsubsection{Bun e ElysiaJS}

O runtime selecionado para os microsserviços foi o \textbf{Bun} \cite{ovensH2023}, uma
plataforma JavaScript/TypeScript de alto desempenho que executa código TypeScript
nativamente, sem etapa de transpilação. Em benchmarks do projeto TechEmpower
\cite{techempower2024}, servidores com Bun processam aproximadamente 2{,}5 a 3 vezes
mais requisições por segundo do que equivalentes em Node.js, tornando-o adequado para
microsserviços em contêineres com recursos limitados.

Sobre o Bun, o framework \textbf{ElysiaJS} \cite{saltyaom2023} oferece validação de
tipos de ponta a ponta. Os esquemas TypeBox definem tanto a validação em tempo de
execução quanto os tipos TypeScript — eliminando divergência entre o contrato da API
e sua implementação. O padrão MVC orientado a domínios é adotado: cada módulo possui
um \textit{controller} (instância Elysia com as rotas), um \textit{service} (lógica de
negócio em classe abstrata) e um \textit{model} (esquemas TypeBox e tipos exportados).

\subsubsection{SQLite via bun:sqlite}

Para a persistência local de cada microsserviço, optou-se pelo \textbf{SQLite} acessado
por meio do driver nativo \texttt{bun:sqlite}, embutido no próprio runtime sem
dependência externa. O SQLite oferece propriedades ACID completas \cite{date2003}. A
opção por \texttt{PRAGMA journal\_mode=WAL} (Write-Ahead Logging) permite múltiplas
leituras concorrentes mesmo durante uma escrita. Para o controle de concorrência,
emprega-se \textbf{versionamento otimista de linhas} \cite{bernstein1987}: a tabela de
assentos possui uma coluna \texttt{version} incrementada a cada atualização. Ao
reservar, a transação verifica se a versão lida ainda é corrente antes de atualizar —
se outro processo modificou o registro no intervalo, o UPDATE retorna \texttt{changes=0}
e a operação é abortada.

\subsubsection{Redis}

O \textbf{Redis} \cite{sanfilippo2009} é empregado em dois papéis. Primeiro, como
\textbf{cache de buscas}: resultados de consultas de voos são armazenados com TTL de 5
minutos (chave \texttt{flights:search:\{origin\}:\{dest\}:\{date\}:\{class\}}), reduzindo
carga sobre o banco nos picos de tráfego. Segundo, como \textbf{mutex distribuído} entre
instâncias do serviço de reservas, detalhado na Seção~2.5.

\subsubsection{RabbitMQ}

A comunicação assíncrona entre serviços é realizada via \textbf{RabbitMQ}
\cite{pivotal2007}, implementando o protocolo AMQP 0-9-1. A topologia adotada é um
\textit{topic exchange} denominado \texttt{airline.events}, no qual mensagens são
roteadas por chaves no formato \texttt{\{agregado\}.\{evento\}} (e.g.,
\texttt{payment.requested}, \texttt{reservation.confirmed}). As filas são duráveis,
garantindo persistência em caso de reinicialização do broker. Cada consumidor utiliza
\texttt{prefetch\_count=1}, e um Dead Letter Exchange (DLX) recebe mensagens que
excedam o número máximo de tentativas, permitindo análise de falhas sem perda de dados.

\subsubsection{NGINX}

O \textbf{NGINX} \cite{reese2008} atua como API Gateway, provendo: roteamento por
prefixo de caminho (\texttt{/api/flights/}, \texttt{/api/reservations/},
\texttt{/api/payments/}); balanceamento de carga pelo algoritmo \texttt{least\_conn}
para serviços com múltiplas réplicas; e rate limiting via \texttt{limit\_req\_zone},
configurado para no máximo 20 requisições por segundo por endereço IP.

\subsubsection{Docker e Docker Compose}

A conteinerização via \textbf{Docker} garante portabilidade e reprodutibilidade do
ambiente. O \textbf{Docker Compose} orquestra as dependências mediante \texttt{depends\_on}
com \texttt{condition: service\_healthy}: o NGINX aguarda todos os microsserviços
saudáveis; os microsserviços aguardam Redis e RabbitMQ passarem nos \textit{health checks}
antes de iniciar. Os dados persistentes são armazenados em volumes nomeados,
sobrevivendo a reinicializações. O sistema utiliza duas redes Docker isoladas:
\texttt{airline-public} (NGINX $\leftrightarrow$ microsserviços) e \texttt{airline-internal}
(microsserviços $\leftrightarrow$ Redis/RabbitMQ).

\subsection{Organização dos Microsserviços}

O sistema é composto por três microsserviços com responsabilidades e fronteiras de dados
bem delimitadas:

\textbf{Serviço de Catálogo de Voos} (\texttt{flight-catalog}, porta 3001): Responsável
pelo inventário de voos — companhias, aeroportos, horários, preços e mapas de assentos.
Por ser o serviço com maior volume de leituras, mantém cache Redis dos resultados de
busca (TTL 5 min). Consome os eventos \texttt{reservation.confirmed} e
\texttt{reservation.cancelled} via RabbitMQ para manter o status dos assentos atualizado.
Banco de dados SQLite com tabelas \texttt{airlines}, \texttt{airports}, \texttt{flights}
e \texttt{seats} (com coluna \texttt{version}).

\textbf{Serviço de Reservas} (\texttt{reservation-service}, porta 3002): Orquestra o
fluxo crítico: aquisição de lock distribuído Redis, verificação e escrita atômica em
SQLite, publicação do evento \texttt{payment.requested}, e atualização do status após
recebimento do resultado de pagamento. Banco SQLite com tabelas \texttt{reservations}
e \texttt{reservation\_history} (trilha de auditoria). Redis key:
\texttt{lock:seat:\{flightId\}:\{seatId\}} com TTL de 30s.

\textbf{Serviço de Pagamentos} (\texttt{payment-service}, porta 3003): Consumidor
assíncrono de \texttt{payment.requested}. Implementa \textbf{idempotência} via Redis
(chave \texttt{payments:idempotency:\{key\}} com TTL de 24h) para evitar processamento
duplicado em caso de reentrega. Após autorização (simulada, 90\% de sucesso), publica
\texttt{payment.succeeded} ou \texttt{payment.failed}. Banco SQLite com tabelas
\texttt{payment\_orders} e \texttt{notifications}.

\subsection{Fluxo de Comunicação}

O fluxo completo de uma reserva ilustra a interação entre todos os componentes:

\begin{enumerate}
  \item O cliente envia \texttt{POST /api/reservations/} ao NGINX (porta 80).
  \item O NGINX encaminha ao \texttt{reservation-service}.
  \item O serviço tenta adquirir o lock Redis: \texttt{SET lock:seat:F1:S22 <token> NX EX 30}.
        Se o lock já estiver tomado, retorna imediatamente \texttt{HTTP 409 Conflict}.
  \item Com o lock adquirido, abre transação SQLite: verifica disponibilidade (lock
        otimista na versão), insere reserva em \texttt{pending}, atualiza assento para
        \texttt{reserved}.
  \item Publica \texttt{payment.requested} no RabbitMQ e retorna \texttt{HTTP 202 Accepted}.
  \item O lock Redis é liberado via script Lua atômico.
  \item O \texttt{payment-service} consome \texttt{q.payment.requested}, processa o
        pagamento (com verificação de idempotência) e publica \texttt{payment.succeeded}
        ou \texttt{payment.failed}.
  \item O \texttt{reservation-service} consome o resultado, atualiza o status da reserva
        (\texttt{confirmed} ou \texttt{cancelled}) e publica o evento correspondente.
  \item O \texttt{flight-catalog} consome o evento e atualiza o status do assento.
\end{enumerate}

A separação entre camada síncrona (REST) e assíncrona (RabbitMQ) garante que o cliente
receba resposta imediata — a latência percebida é da etapa de lock e transação, não do
processamento de pagamento.

\subsection{Gerenciamento de Concorrência}

\subsubsection{O Problema do Double-Booking}

Em sistemas de reserva com alta concorrência, dois usuários podem enviar
simultaneamente requisições para o mesmo assento. Na ausência de coordenação, ambos
leriam o assento como \texttt{available}, passariam pela validação e concluiriam a
escrita — resultando em duas reservas válidas para um único assento
\cite{gray92}. Este é um caso clássico de \textit{lost update}, anomalia em que
atualizações concorrentes se sobrescrevem silenciosamente.

\subsubsection{Lock Distribuído com Redis}

A solução utiliza o Redis como mutex distribuído, baseando-se na atomicidade do comando
\texttt{SET key value NX EX ttl} \cite{antirez2016}. O sufixo \texttt{NX}
(\textit{Not eXists}) instrui o Redis a realizar a atribuição somente se a chave não
existir; \texttt{EX ttl} define expiração automática. Como o Redis é
\textit{single-threaded}, esta operação é garantidamente atômica — não há condição de
corrida:

\begin{lstlisting}[language={}]
  Processo A: SET lock:seat:F1:S22 "token-a" NX EX 30  -> OK   (lock adquirido)
  Processo B: SET lock:seat:F1:S22 "token-b" NX EX 30  -> nil  (lock recusado)
\end{lstlisting}

A liberação do lock é realizada via script Lua atômico, garantindo que apenas o
proprietário do lock possa liberá-lo:

\begin{lstlisting}[language={}]
  if redis.call("GET", key) == token then
    redis.call("DEL", key)
  end
\end{lstlisting}

\subsubsection{Versionamento Otimista como Segunda Camada}

O lock Redis é a defesa primária, mas o versionamento otimista no SQLite serve como
segunda camada de segurança para casos extremos (e.g., expiração do TTL durante seção
crítica longa). O UPDATE só persiste se a versão do assento não mudou desde o SELECT:

\begin{lstlisting}[language=SQL]
  UPDATE seats SET status = 'reserved', version = version + 1
  WHERE id = ? AND version = ?  -- falha se version mudou
\end{lstlisting}

Se \texttt{changes=0}, a transação é abortada com \texttt{HTTP 409}.

\subsubsection{Mecanismo Anti-Deadlock}

O TTL de 30 segundos no lock Redis garante que, em caso de falha do processo durante a
seção crítica (crash ou desconexão), o lock seja automaticamente liberado sem intervenção
manual. Reservas em status \texttt{pending} além do prazo \texttt{expires\_at} (15 min)
são elegíveis para cancelamento automático, com liberação do assento.

A combinação Redis + SQLite forma uma defesa em profundidade:

\begin{table}[ht]
\centering
\caption{Camadas de proteção contra double-booking}
\label{tab:layers}
\begin{tabular}{|l|l|l|}
\hline
\textbf{Camada} & \textbf{Mecanismo} & \textbf{Protege contra} \\
\hline
Aplicação & Redis \texttt{SET NX EX} (mutex) & Requisições concorrentes ao mesmo assento \\
\hline
Banco de dados & \texttt{version} otimista + \texttt{changes} & Atualizações simultâneas na mesma linha \\
\hline
\end{tabular}
\end{table}

\subsection{Registro de Eventos e Visualização}

Para fins de observabilidade e validação, o \texttt{reservation-service} registra todos
os eventos relevantes em uma tabela SQLite \texttt{events} (auto-limitada a 500
entradas por trigger). Um endpoint \texttt{GET /api/events?since=\{id\}\&limit=\{n\}}
expõe esses eventos, consumidos pelo painel de arquitetura em tempo real implementado
na página \texttt{architecture.html}. O painel utiliza polling a cada 1,2 segundos e
anima graficamente as interações entre nós (NGINX, serviços, Redis, RabbitMQ),
distinguindo eventos novos de eventos históricos por meio de uma flag de inicialização.


% =============================================================================
\section{Resultados e Discussão}
% =============================================================================

\subsection{Comportamento sob Carga Concorrente}

Para validar a prevenção de double-booking, foram executados testes manuais com
múltiplas requisições concorrentes utilizando a ferramenta \texttt{curl} em paralelo e
o próprio cliente web do sistema. O cenário de teste consistiu em enviar
simultaneamente $N$ requisições \texttt{POST /api/reservations/} para o mesmo assento
com diferentes identificadores de usuário.

Os resultados confirmaram o comportamento esperado: apenas a primeira requisição a
adquirir o lock Redis obteve sucesso (\texttt{HTTP 202 Accepted}); todas as demais
receberam \texttt{HTTP 409 Conflict} com a mensagem
\texttt{seat\_temporarily\_locked}, sem jamais acessar o banco de dados. Este
comportamento é fundamental: o Redis absorve a contenção antes de qualquer I/O no
SQLite, mantendo a latência de rejeição inferior a 5 ms.

\subsection{Latência do Pipeline de Reserva}

O pipeline completo de uma reserva — do \texttt{POST} inicial até a atualização do
status para \texttt{confirmed} — é composto por duas fases com latências distintas:

\begin{itemize}
  \item \textbf{Fase síncrona} (lock Redis + transação SQLite + publicação RabbitMQ):
        aproximadamente 80--120 ms, determinando a latência percebida pelo cliente.
  \item \textbf{Fase assíncrona} (processamento do pagamento pelo \texttt{payment-service}
        + atualização do status): aproximadamente 200--400 ms adicionais, transparentes
        ao cliente (que recebeu \texttt{202 Accepted}).
\end{itemize}

O uso do padrão assíncrono via RabbitMQ garante que o cliente não aguarde o
processamento do pagamento — padrão fundamental em sistemas de e-commerce de alta
disponibilidade \cite{kleppmann2017}.

\subsection{Tolerância a Falhas e Idempotência}

A idempotência no \texttt{payment-service} foi validada reprocessando manualmente
mensagens da fila (via RabbitMQ Management UI em \texttt{localhost:15672}). Em todos
os casos, a segunda tentativa de processar o mesmo \texttt{idempotency\_key} foi
descartada silenciosamente, sem duplicação de registros de pagamento.

O Dead Letter Exchange demonstrou seu funcionamento ao simular uma falha no processamento
(injeção de payload inválido): após 3 tentativas com \texttt{nack}, a mensagem foi
encaminhada para a fila DLX sem perda, podendo ser inspecionada e reprocessada
manualmente.

\subsection{Cancelamento e Consistência de Dados}

O cancelamento de reservas implementa a restauração do assento em duas frentes: (i) na
operação de cancelamento explícito (\texttt{DELETE /api/reservations/:id}), onde a
transação SQLite atualiza simultaneamente o status da reserva e o assento local; e (ii)
no consumidor de \texttt{payment.failed}, que invoca \texttt{restoreSeat()} para
devolver o assento ao estado \texttt{available}. Esta redundância garante que o assento
sempre seja liberado, independentemente do caminho de falha.

\subsection{Painel de Arquitetura em Tempo Real}

O painel desenvolvido em \texttt{architecture.html} utiliza SVG puro com animações CSS
para visualizar as interações entre serviços em tempo real. Cada evento registrado
(aquisição de lock, transação, publicação de mensagem, confirmação de pagamento) é
representado graficamente como uma partícula animada percorrendo a aresta correspondente
no grafo de arquitetura. O painel diferencia eventos históricos (carregados silenciosamente
na inicialização) de eventos novos (animados), eliminando animações espúrias em
atualizações de página.

\subsection{Limitações e Trabalhos Futuros}

A principal limitação da implementação atual é o uso de SQLite como banco de dados
local: o arquivo \texttt{.db} não suporta múltiplos writers em contêineres distintos
que compartilhem o mesmo volume. Portanto, ao escalar o \texttt{reservation-service}
horizontalmente (múltiplas réplicas), é necessário garantir que cada réplica acesse
seu próprio arquivo SQLite — o que torna o lock Redis ainda mais crítico como único
ponto de coordenação distribuída. Em produção, o caminho recomendado seria migrar
para PostgreSQL com particionamento de dados, mantendo o mecanismo de lock Redis.

Outras limitações incluem: (i) o processamento de pagamento é simulado (mock), sem
integração com gateway real; (ii) não há mecanismo de compensação automática para
reservas expiradas em \texttt{pending}; e (iii) a autenticação de usuários não está
implementada, sendo o \texttt{user\_id} informado livremente pelo cliente.

Como trabalhos futuros, identificam-se: implementação de Saga Pattern para transações
distribuídas entre serviços; integração com gateway de pagamento real; autenticação JWT
via NGINX; e substituição do SQLite por PostgreSQL com suporte a múltiplas réplicas
com escrita concorrente.


% =============================================================================
\section{Conclusão}
% =============================================================================

Este trabalho apresentou o AVSYS, um sistema distribuído de reserva de passagens aéreas
que aplica na prática os principais conceitos estudados na disciplina de Sistemas
Distribuídos e Programação Paralela. A arquitetura de microsserviços com comunicação
assíncrona via RabbitMQ demonstrou a viabilidade de construir sistemas desacoplados e
tolerantes a falhas sem recorrer a transações distribuídas globais.

O principal contribuição técnica é a estratégia de prevenção de double-booking em duas
camadas: o lock distribuído Redis elimina a contenção antes do acesso ao banco de dados,
enquanto o versionamento otimista SQLite serve como segunda barreira de segurança. Esta
combinação garante corretude sem sacrificar a disponibilidade do sistema — requisição
rejeitada por contenção recebe resposta em menos de 5 ms, sem bloquear outras
operações.

Os resultados experimentais validaram os mecanismos de concorrência, idempotência e
tolerância a falhas, e o painel de arquitetura em tempo real tornou os fluxos
distribuídos observáveis e compreensíveis. O sistema é executável em qualquer ambiente
com Docker Compose disponível, com configuração zero adicional.

\bibliographystyle{sbc}

\begin{thebibliography}{99}

\bibitem{antirez2016}
ANTIREZ, S. \textbf{Distributed locks with Redis}. Redis Blog, 2016.
Disponível em: https://redis.io/docs/manual/patterns/distributed-locks/

\bibitem{bernstein1987}
BERNSTEIN, P. A.; HADZILACOS, V.; GOODMAN, N. \textbf{Concurrency Control and
Recovery in Database Systems}. Addison-Wesley, 1987.

\bibitem{brewer2000}
BREWER, E. A. Towards robust distributed systems. In: \textbf{Proceedings of the
19th Annual ACM Symposium on Principles of Distributed Computing (PODC)}, p.~7,
Portland, OR, USA, 2000.

\bibitem{date2003}
DATE, C. J. \textbf{An Introduction to Database Systems}. 8. ed. Pearson, 2003.

\bibitem{gray92}
GRAY, J.; REUTER, A. \textbf{Transaction Processing: Concepts and Techniques}.
Morgan Kaufmann, 1992.

\bibitem{kleppmann2017}
KLEPPMANN, M. \textbf{Designing Data-Intensive Applications}. O'Reilly Media, 2017.

\bibitem{newman2015}
NEWMAN, S. \textbf{Building Microservices}. O'Reilly Media, 2015.

\bibitem{ovensH2023}
OVEN SH. \textbf{Bun — A fast all-in-one JavaScript runtime}. 2023.
Disponível em: https://bun.sh

\bibitem{pivotal2007}
PIVOTAL SOFTWARE. \textbf{RabbitMQ — Messaging that just works}. 2007.
Disponível em: https://www.rabbitmq.com

\bibitem{reese2008}
REESE, W. Nginx: the High-Performance Web Server and Reverse Proxy.
\textbf{Linux Journal}, n.~173, 2008.

\bibitem{richardson2018}
RICHARDSON, C. \textbf{Microservices Patterns}. Manning Publications, 2018.

\bibitem{saltyaom2023}
SALTYAOM. \textbf{ElysiaJS — Ergonomic Framework for Humans}. 2023.
Disponível em: https://elysiajs.com

\bibitem{sanfilippo2009}
SANFILIPPO, S. \textbf{Redis — Remote Dictionary Server}. 2009.
Disponível em: https://redis.io

\bibitem{techempower2024}
TECHEMPOWER. \textbf{TechEmpower Framework Benchmarks}. 2024.
Disponível em: https://www.techempower.com/benchmarks/

\end{thebibliography}

\end{document}
